% 3. ML-обнаружение
\begin{frame}{ML-обнаружение вредоносного трафика}
\framesubtitle{Новая парадигма сетевой безопасности}
\vspace{1mm}
\begin{columns}[T,onlytextwidth]
\column{0.44\textwidth}
\begin{card}[height=5.9cm]
  \cardtitle{Почему ML вытесняет правила}\vspace{1mm}
  \iconitem{steel}{ic_eye_w.png}{Работает с зашифрованным трафиком — не смотрит в payload, только статистика}\vspace{4mm}
  \iconitem{steel}{ic_bolt_w.png}{Ловит неизвестные и сложные атаки лучше сигнатурных методов}\vspace{4mm}
  \iconitem{steel}{ic_server_w.png}{Развёртывается на «железе»: программируемые коммутаторы, SmartNIC}
\end{card}
\column{0.53\textwidth}
\begin{wcard}[height=5.9cm]
  \cardtitle{Признаки, на которых учатся детекторы}
  \centering\includegraphics[width=0.94\linewidth]{fig4.png}\par\vspace{1.2mm}
  \raggedright\figcap{Размеры пакетов и межпакетные задержки (IPD) — базовые статистики потока}\par\vspace{2mm}
  \begin{darkcard}[top=1.8mm,bottom=1.8mm]\scriptsize
    Примеры детекторов: Whisper, FlowLens, NetBeacon, Kitsune, \mbox{CICFlowMeter+MLP}
  \end{darkcard}
\end{wcard}
\end{columns}
\end{frame}
